\documentclass[11pt]{article}
\usepackage[T1]{fontenc}
\usepackage{amsmath,amssymb,amsthm,mathtools}
\usepackage[margin=1.1in]{geometry}
\usepackage{hyperref}
\newtheorem{theorem}{Theorem}[section]
\newtheorem{lemma}[theorem]{Lemma}
\newtheorem{corollary}[theorem]{Corollary}
\newcommand{\I}{\mathcal I}
\newcommand{\R}{\mathbb R}
\title{An explicit constrained exponential pair problem}
\author{Ziyang Qin}
\date{}
\begin{document}
\maketitle

\begin{abstract}
We determine the minimum of an exponential interaction over continuous
functions with prescribed lower and upper obstacles. The minimizer consists
of two constant pieces joined by an explicit free arc. Its value is given
for every cap, and an exact convex deficit identity proves uniqueness.
At every nondegenerate cap the minimizer has derivative jumps at both
contacts; every admissible function differentiable in the interior has
strictly larger value.
\end{abstract}

\section{The constrained problem}

For $C\geq1$, consider real functions whose restrictions to the unit interval
are continuous, and let
\[
 \mathcal A_C=\{k:\R\to\R:k|_{[0,1]}\in C([0,1]),\quad
      0\leq k(v)\leq\log C\text{ for }v\in[0,1]\},
\]
\[
 \I(k)=\int_0^1\int_0^s 2s e^{k(v)-k(s)}\,dv\,ds.
\]
All functions and integrals in this paper are real-valued. The interval
constraints include both endpoints. In particular $\mathcal A_1$ contains
only functions vanishing on $[0,1]$. Values outside this interval do not
affect the functional; uniqueness always means equality on $[0,1]$.

\begin{theorem}\label{thm:minimum}
For every $C\geq1$ there is a unique $q\geq1$ satisfying
\begin{equation}\label{eq:cap}
 \log q+\frac23(q^3-1)=\log C.
\end{equation}
Set
\begin{equation}\label{eq:contacts}
 \alpha=\frac1{\sqrt{q^4+2q}},\qquad \beta=\alpha q^2.
\end{equation}
Then $0<\alpha\leq\beta<1$, and the function $k_*=\log E_*$, where
\begin{equation}\label{eq:candidate}
 E_*(v)=
 \begin{cases}
 1,&v\leq\alpha,\\
 \displaystyle\sqrt{v/\alpha}\,
  \exp\!\left(\frac23\bigl((v/\alpha)^{3/2}-1\bigr)\right),
   &\alpha\leq v\leq\beta,\\
 C,&v\geq\beta,
 \end{cases}
\end{equation}
belongs to $\mathcal A_C$. For every $k\in\mathcal A_C$,
\begin{equation}\label{eq:minimum}
 \I(k)\geq\frac23-\beta+\frac1{3\beta},
\end{equation}
with equality if and only if $k=k_*$ on $[0,1]$.
At $C=1$ the profile is identically zero and the minimum is $2/3$.
\end{theorem}

The degeneracy in the middle interval is essential at cap one. No derivative
formula on a nonempty free interval is needed for that endpoint.

\begin{lemma}\label{lem:parameter}
The parameter in \eqref{eq:cap} exists uniquely and increases strictly with
$C\geq1$. It equals one precisely when $C=1$. The contacts satisfy
\begin{equation}\label{eq:matching}
 \beta^2=\frac{q^3}{q^3+2},\qquad
 1-\beta^2=2\alpha^2q,\qquad
 \frac\alpha q=\frac{1-\beta^2}{2\beta}.
\end{equation}
The contacts are distinct precisely when $C>1$.
\end{lemma}

\begin{proof}
The function $F(q)=\log q+\frac23(q^3-1)$ is continuous and strictly
increasing on $[1,\infty)$. It has $F(1)=0$ and $F(C)\geq\log C$
for $C\geq1$, so the intermediate value theorem on $[1,C]$ gives the
unique parameter. Strict increase of $F$ and of the logarithm gives the
claimed dependence on $C$. Squaring \eqref{eq:contacts} gives
\eqref{eq:matching}. In particular $\beta<1$, while $q\geq1$ gives
$\alpha\leq\beta$, strictly when $q>1$.
\end{proof}

The free expression in \eqref{eq:candidate} equals one at $v=\alpha$
and $q\exp(\frac23(q^3-1))=C$ at $v=\beta$. Both its square-root
and exponential factors are nonnegative and increasing for positive $v$.
It therefore joins the obstacles continuously and stays between their
values. This proves the asserted admissibility.

\section{Exact primitives and marginal}

Write $\phi(v)=\frac23((v/\alpha)^{3/2}-1)$ on the free arc, and define
the actual prefix and tail integrals
\[
 P(v)=\int_0^v E_*(t)\,dt,\qquad
 Q(v)=\int_v^1\frac{2s}{E_*(s)}\,ds.
\]
\begin{lemma}\label{lem:primitives}
The two integrals have the formulas
\begin{equation}\label{eq:primitives}
 P(v)=
 \begin{cases}
 v,&v\leq\alpha,\\
 \alpha e^{\phi(v)},&\alpha\leq v\leq\beta,\\
 C(v-\beta+\alpha/q),&v\geq\beta,
 \end{cases}
 \quad
 Q(v)=
 \begin{cases}
 3\alpha^2-v^2,&v\leq\alpha,\\
 2\alpha^2 e^{-\phi(v)},&\alpha\leq v\leq\beta,\\
 (1-v^2)/C,&v\geq\beta,
 \end{cases}
\end{equation}
for $v\in[0,1]$. Consequently the marginal density
$A(v)=E_*(v)Q(v)-2vP(v)/E_*(v)$ is
\begin{equation}\label{eq:marginal}
 A(v)=
 \begin{cases}
 3(\alpha^2-v^2),&0\leq v\leq\alpha,\\
 0,&\alpha\leq v\leq\beta,\\
 -3(v-\beta)\bigl(v+1/(3\beta)\bigr),&\beta\leq v\leq1.
 \end{cases}
\end{equation}
\end{lemma}

\begin{proof}
Direct differentiation on the positive half-line gives
\[
 \frac{d}{dv}\bigl(\alpha e^{\phi(v)}\bigr)=E_*(v),\qquad
 \frac{d}{dv}\bigl(-2\alpha^2 e^{-\phi(v)}\bigr)
 =\frac{2v}{E_*(v)}
\]
on the free arc. At its upper endpoint the first primitive has value
$C\alpha/q$, while the second has value $-(1-\beta^2)/C$ by
\eqref{eq:matching}. Splitting the integrals at the contacts gives
\eqref{eq:primitives}; the calculation also applies when the contacts
coincide. Substitution into the definition of $A$ proves the first two
pieces of \eqref{eq:marginal}. On the upper piece it gives
$1-3v^2+2v\beta-2v\alpha/q$, which has the stated factorization by
the last identity in \eqref{eq:matching}.
\end{proof}

\section{Convex calibration and rigidity}

For continuous $k,g$ on $[0,1]$, put $\delta=k-g$ and
\[
 A_g(v)=e^{g(v)}\int_v^1 2s e^{-g(s)}\,ds
       -2v e^{-g(v)}\int_0^v e^{g(t)}\,dt.
\]
\begin{lemma}\label{lem:deficit}
The exact identity
\begin{align}\label{eq:deficit}
 \I(k)-\I(g)
 &=\int_0^1 A_g(v)\delta(v)\,dv+\mathcal R(k,g),\\
 \mathcal R(k,g)
 &=\int_0^1\int_0^s 2s e^{g(v)-g(s)}
 \left[e^{\delta(v)-\delta(s)}-1-\delta(v)+\delta(s)\right]dv\,ds
 \nonumber
\end{align}
holds. The remainder is nonnegative, and it vanishes if and only if
$\delta$ is constant on $[0,1]$.
\end{lemma}

\begin{proof}
Apply the identity $e^{x+y}-e^x=e^xy+e^x(e^y-1-y)$ with
$x=g(v)-g(s)$ and $y=\delta(v)-\delta(s)$. Fubini's theorem identifies
the linear term with the stated marginal integral. All integrands are
continuous on the compact square before restriction to the triangle, so
these integrations are justified. The inequality $e^y\geq1+y$ proves
nonnegativity. If the remainder vanishes, its nonnegative integrand is
zero almost everywhere. Continuity makes it zero at every point of the
open triangle $0<v<s<1$: otherwise it would be positive on an open set
of positive measure. The prefactor is positive there, and strictness of
$e^y>1+y$ for $y\ne0$ gives $\delta(v)=\delta(s)$. Any two interior
points can be ordered, so $\delta$ is constant in the interior, and
continuity extends that constant to the endpoints. The converse follows
directly from the integrand.
\end{proof}

\begin{proof}[Proof of Theorem~\ref{thm:minimum}]
Take $g=k_*$. On the lower obstacle $\delta=k\geq0$ and $A\geq0$;
on the upper obstacle $\delta=k-\log C\leq0$ and $A\leq0$; on the
free arc $A=0$. Thus both terms in \eqref{eq:deficit} are nonnegative,
and $k_*$ is a minimizer. If equality holds, Lemma~\ref{lem:deficit}
gives $k-k_*=c$ on the interval. At zero the lower constraint gives
$c\geq0$, while at one the upper constraint gives $c\leq0$. Therefore
$c=0$.

It remains to evaluate the value. Fubini gives
$\I(k_*)=\int_0^1 2vP(v)/E_*(v)\,dv$. The integrand on the three
pieces is, respectively,
\[
 2v^2,\qquad 2\alpha^2\sqrt{v/\alpha},\qquad
 2v(v-\beta+\alpha/q).
\]
Its integral is consequently
\[
 \frac23\alpha^3+\frac43\alpha^3(q^3-1)
 +\frac23(1-\beta^3)-(\beta-\alpha/q)(1-\beta^2).
\]
The identities \eqref{eq:matching} reduce this expression to
$\frac23-\beta+1/(3\beta)$. The integration of the middle term is
valid also on a degenerate interval. Alternatively, at $C=1$ one
directly has $\I(0)=\int_0^1 2s^2\,ds=2/3$.
\end{proof}

\begin{theorem}\label{thm:contacts}
For $C>1$, the one-sided derivatives of the minimizer are
\[
 k'_{*,-}(\alpha)=0,\qquad k'_{*,+}(\alpha)=\frac3{2\alpha},
 \qquad
 k'_{*,-}(\beta)=\frac1{2\beta}+\frac q\alpha,
 \qquad k'_{*,+}(\beta)=0.
\]
In particular, every $k\in\mathcal A_C$ differentiable on $(0,1)$
satisfies strict inequality in \eqref{eq:minimum}.
\end{theorem}

\begin{proof}
The obstacles are constant. In the free arc direct differentiation gives
$k_*'(v)=1/(2v)+\sqrt{v/\alpha}/\alpha$. Since $\alpha<\beta$, its
one-sided limits give the displayed positive derivatives. Both contacts
are interior, so the jumps exclude differentiability there. Equality
in \eqref{eq:minimum} would identify a differentiable competitor with
$k_*$ throughout the interval, contradicting either jump.
\end{proof}

\section{A quantitative exponential estimate}

\begin{lemma}\label{lem:quadratic}
For real $m,x,y$ with $m\leq x,y$,
\[
 e^x-e^y-e^y(x-y)\geq\frac{e^m}{2}(x-y)^2.
\]
\end{lemma}
\begin{proof}
The function $t\mapsto e^t-e^mt^2/2$ has nonnegative second derivative
on $[m,\infty)$. Its supporting-line inequality at $y$, applied to $x$,
is exactly the stated bound after rearrangement. This applies in either
order of $x$ and $y$, including the boundary value $m$.
\end{proof}

\end{document}
